% Flavor catalog per-process draft.
% process_id: L004
% family: charged_lepton
% owner_agent_id: PKA-L004
% writer_agent_id: null
% checker_agent_id: null
% opus_signoff_id: null
% Canonical metadata sidecar: L004.yaml
% Status is controlled by the YAML sidecar status_history, not this comment.

\section{L004: \texorpdfstring{\(\mu-e\)}{mu-e} Conversion in Gold}

\subsection*{Process}
\textbf{Standard notation:} \(R_{\mu e}^{\rm Au} =
\Gamma(\mu^-{\rm Au}\to e^-{\rm Au})/
\Gamma_{\rm capture}(\mu^-{\rm Au})\).
This entry covers coherent neutrinoless muon-to-electron conversion in a
gold target, the strongest published SINDRUM II conversion bound and the
direct companion to the aluminum and titanium conversion entries.

\subsection*{PDG-or-equivalent value}
The streamed PDG live muon listing, generated on 2026-05-14 and accessed on
2026-05-16, gives
\[
  R_{\mu e}^{\rm Au} < 7\times 10^{-13}\qquad (90\%\ {\rm C.L.}) .
\]
These dates, the \(7\times10^{-13}\) value, and the \(90\%\) C.L. qualifier are
recorded in
\texttt{L004.yaml:pdg\_or\_equivalent.primary\_current\_limit}.  The PDG line
identifies the document as BERTL 06, technique SPEC, comment SINDRUM II.  The
local minimal text snapshot is
\texttt{flavor\_catalog/references/L004/pdg2026\_muon\_listing\_s004\_au\_conversion.txt}.
The SINDRUM II article metadata is separately snapshotted from Crossref for
\textit{Eur. Phys. J. C} \textbf{47}, 337--346 (2006), DOI
\texttt{10.1140/epjc/s2006-02582-x}
(\texttt{L004.yaml:pdg\_or\_equivalent.experimental\_publication}).

\subsection*{Relevance to RS / anarchic flavor}
In a warped lepton-flavor extension, \(\mu-e\) conversion probes more than the
dipole structure tested by \(\mu\to e\gamma\).  Tree-level or loop-induced
lepton-quark vector and scalar operators, flavor-changing \(Z\)-like
couplings, Higgs-mediated scalar currents, and dipoles can all contribute to
the same coherent nuclear final state.  A heavy target such as gold is useful
as an existing benchmark, but its interpretation is target and operator
dependent.

\subsection*{Post-2008 developments}
Since the arXiv:0804.1954 RS-flavor baseline, no newer gold-target result has
replaced the SINDRUM II bound in the PDG listing
(\texttt{L004.yaml:pdg\_or\_equivalent.primary\_current\_limit};
\texttt{L004.yaml:pdg\_or\_equivalent.post\_2008\_baseline}).  The major
change is the experimental program around lighter targets: the Mu2e Technical
Design Report
states an aluminum conversion search aiming for sensitivity approximately four
orders of magnitude beyond the previous world-best conversion limits
(\texttt{L004.yaml:pdg\_or\_equivalent.post\_2008\_developments[0]}), and the
COMET Phase-I TDR quotes an aluminum sensitivity goal \(3.1\times10^{-15}\)
with an expected \(90\%\) upper limit \(7.0\times10^{-15}\)
(\texttt{L004.yaml:pdg\_or\_equivalent.post\_2008\_developments[1]}).  Thus the
gold limit remains the published benchmark, while future comparisons will need
target-by-target translations.

\subsection*{Constraint validity / model dependence}
The experimental statement is robust: it is a null search normalized to the
ordinary muon-capture rate on gold.  The mapping to RS parameters is not a
single-number rescaling.  It requires a lepton-quark effective operator basis,
gold nuclear overlap integrals and capture normalization, and a convention for
combining dipole, scalar, and vector amplitudes.  Comparisons with aluminum or
titanium limits are therefore model- and operator-dependent.

\subsection*{Code coverage in this repo}
\textbf{NO.}  The required greps over \texttt{quarkConstraints/}, \texttt{qcd/},
\texttt{flavorConstraints/}, \texttt{neutrinos/}, \texttt{yukawa/},
\texttt{warpConfig/}, \texttt{solvers/}, \texttt{scanParams/}, and
\texttt{tests/} found no SINDRUM, Mu2e, COMET, \(R_{\mu e}\), gold conversion,
or muon-to-electron conversion implementation.  The only adjacent charged-LFV
code is the dipole-only \(\mu\to e\gamma\) checker at
\texttt{flavorConstraints/muToEGamma.py}:75, called by the scan at
\texttt{scanParams/scan.py}:524.

\subsection*{Implementation difficulty}
\textbf{HIGH.}  A production constraint needs a new \(\mu-e\) conversion
observable, lepton-quark Wilson coefficients, target-specific gold nuclear
inputs, capture-rate normalization, and likely shared EFT matching/running
with \(\mu\to e\gamma\) and \(\mu\to3e\).  This is a new mode calculation, not
an input-value update to the existing dipole checker.

\subsection*{Key references}
Process-local source keys before bibliography consolidation:
\texttt{PDG2026\_MuonAuConversion}, \texttt{SINDRUMII2006\_GoldMuE},
\texttt{Mu2eTDR2015}, \texttt{COMETPhaseITDR2020}, and \texttt{CFW2008}.
